% GENERATED FILE. Edit canonical lessons, then run npm run generate:books.
% canonical-source-hash: fnv1a64:33dc6358434dc6b0
% canonical-lessons: GE-C382-verbot, GE-C382-zulaessig, GE-C382-strafbar, GE-C382-verstossen, GE-C382-verstoss

\chapter{Gegen die Regeln}
\label{ch:ge-gegen-die-regeln}
\hlchaptermodality{382}

\begin{chapteropening}
\textbf{By the end of this chapter:} \emph{I can talk about bans and what is permitted, about offences, and about breaking a rule.}

Complete the last lesson of chapter 382: I can talk about bans and what is permitted, about offences, and about breaking a rule.
\end{chapteropening}

\section[das Verbot]{das Verbot — the ban}

\label{lesson:GE-C382-verbot}



\begin{quote}
Before the new one: say the German for to keep to, then the German for obliged.
\end{quote}

\subsection*{You'll want to know: das Verbot}
\textbf{das Verbot} — \textquotedblleft{}the ban\textquotedblright{}.

\begin{quote}
In vielen Parks gibt es ein Verbot für Hunde.
\end{quote}

\begin{quote}
\emph{In many parks there is a ban on dogs.}
\end{quote}

Plural: \emph{die Verbote}. From \emph{verbieten} — to forbid.

\begin{cousinweb}
\emph{Verbot} is the noun of \emph{verbieten} \textquotedblleft{}to forbid\textquotedblright{}, built with \emph{ver-} on \emph{bieten} \textquotedblleft{}to offer, to announce\textquotedblright{}; EtymDB lists English \emph{bid} as a cousin of \emph{bieten}.
\end{cousinweb}

\begin{grammarlens}[title={What you've built}]
One more word for telling, explaining and describing at B1.
\end{grammarlens}

\subsection*{Your turn}
\begin{itemize}
\raggedright
  \item \emph{Say it:} \emph{das Verbot}
  \item \emph{Say it:} \emph{das Verbot}, once more
  \item \emph{Say it:} a sentence of your own with \emph{das Verbot}
  \item \emph{Recall:} say the German for to keep to, then the German for obliged, then say \emph{das Verbot} again
\end{itemize}

\begin{tcolorbox}[breakable,colback=teal!4,colframe=teal!35!black,title={Before you move on}]
What does das Verbot mean? (\textquotedblleft{}The ban\textquotedblright{}.) Say it once more.
\end{tcolorbox}
\section[zulässig]{zulässig — permitted, allowed}

\label{lesson:GE-C382-zulaessig}



\begin{quote}
Before the new one: say the German for obliged, then the German for the ban.
\end{quote}

\subsection*{You'll want to know: zulässig}
\textbf{zulässig} — \textquotedblleft{}permitted, allowed\textquotedblright{}.

\begin{quote}
Auf dieser Straße sind nur dreißig Kilometer pro Stunde zulässig.
\end{quote}

\begin{quote}
\emph{On this road only thirty kilometres an hour are permitted.}
\end{quote}

Formal. The opposite is \emph{unzulässig}.

\begin{cousinweb}
\emph{zulässig} is built with \emph{-ig} on \emph{zulassen} \textquotedblleft{}to allow\textquotedblright{}, from \emph{lassen} \textquotedblleft{}to let\textquotedblright{}, Old High German \emph{lāzan} (EtymDB): what may be let through.
\end{cousinweb}

\begin{grammarlens}[title={What you've built}]
One more word for telling, explaining and describing at B1.
\end{grammarlens}

\subsection*{Your turn}
\begin{itemize}
\raggedright
  \item \emph{Say it:} \emph{zulässig}
  \item \emph{Say it:} \emph{zulässig}, once more
  \item \emph{Say it:} a sentence of your own with \emph{zulässig}
  \item \emph{Recall:} say the German for obliged, then the German for the ban, then say \emph{zulässig} again
\end{itemize}

\begin{tcolorbox}[breakable,colback=teal!4,colframe=teal!35!black,title={Before you move on}]
What does zulässig mean? (\textquotedblleft{}Permitted, allowed\textquotedblright{}.) Say it once more.
\end{tcolorbox}
\section[strafbar]{strafbar — punishable, a criminal offence}

\label{lesson:GE-C382-strafbar}



\begin{quote}
Before the new one: say the German for the ban, then the German for permitted.
\end{quote}

\subsection*{You'll want to know: strafbar}
\textbf{strafbar} — \textquotedblleft{}punishable, a criminal offence\textquotedblright{}.

\begin{quote}
Fahren ohne Führerschein ist strafbar.
\end{quote}

\begin{quote}
\emph{Driving without a licence is a criminal offence.}
\end{quote}

\emph{Sich strafbar machen} — to commit an offence.

\begin{cousinweb}
\emph{strafbar} is built with \emph{-bar} \textquotedblleft{}able to be\textquotedblright{} on \emph{Strafe} \textquotedblleft{}punishment\textquotedblright{}, inherited from Middle High German \emph{strāfe} (EtymDB).
\end{cousinweb}

\begin{grammarlens}[title={What you've built}]
One more word for telling, explaining and describing at B1.
\end{grammarlens}

\subsection*{Your turn}
\begin{itemize}
\raggedright
  \item \emph{Say it:} \emph{strafbar}
  \item \emph{Say it:} \emph{strafbar}, once more
  \item \emph{Say it:} a sentence of your own with \emph{strafbar}
  \item \emph{Recall:} say the German for the ban, then the German for permitted, then say \emph{strafbar} again
\end{itemize}

\begin{tcolorbox}[breakable,colback=teal!4,colframe=teal!35!black,title={Before you move on}]
What does strafbar mean? (\textquotedblleft{}Punishable, a criminal offence\textquotedblright{}.) Say it once more.
\end{tcolorbox}
\section[verstoßen]{verstoßen — to break (a rule), to violate}

\label{lesson:GE-C382-verstossen}



\begin{quote}
Before the new one: say the German for permitted, then the German for punishable.
\end{quote}

\subsection*{You'll want to know: verstoßen}
\textbf{verstoßen} — \textquotedblleft{}to break (a rule), to violate\textquotedblright{}.

\begin{quote}
Der Spieler hat gegen die Regeln verstoßen.
\end{quote}

\begin{quote}
\emph{The player broke the rules.}
\end{quote}

\emph{Gegen etwas verstoßen} — to break a rule. Strong: \emph{er verstößt, er hat verstoßen}.

\begin{cousinweb}
\emph{verstoßen} is built with \emph{ver-} on \emph{stoßen} \textquotedblleft{}to push, to thrust\textquotedblright{}, inherited from Old High German \emph{stōzan} and Proto-Germanic \emph{stautaną} (EtymDB): you push against a rule.
\end{cousinweb}

\begin{grammarlens}[title={What you've built}]
One more word for telling, explaining and describing at B1.
\end{grammarlens}

\subsection*{Your turn}
\begin{itemize}
\raggedright
  \item \emph{Say it:} \emph{verstoßen}
  \item \emph{Say it:} \emph{verstoßen}, once more
  \item \emph{Say it:} a sentence of your own with \emph{verstoßen}
  \item \emph{Recall:} say the German for permitted, then the German for punishable, then say \emph{verstoßen} again
\end{itemize}

\begin{tcolorbox}[breakable,colback=teal!4,colframe=teal!35!black,title={Before you move on}]
What does verstoßen mean? (\textquotedblleft{}To break (a rule), to violate\textquotedblright{}.) Say it once more.
\end{tcolorbox}
\section[der Verstoß]{der Verstoß — the violation, the breach}

\label{lesson:GE-C382-verstoss}



\begin{quote}
Before the new one: say the German for punishable, then the German for to break.
\end{quote}

\subsection*{You'll want to know: der Verstoß}
\textbf{der Verstoß} — \textquotedblleft{}the violation, the breach\textquotedblright{}.

\begin{quote}
Für jeden Verstoß gibt es eine Strafe.
\end{quote}

\begin{quote}
\emph{There is a penalty for every violation.}
\end{quote}

Plural: \emph{die Verstöße}. \emph{Ein Verstoß gegen} + accusative.

\begin{cousinweb}
\emph{Verstoß} is the noun of \emph{verstoßen}, built on \emph{stoßen} \textquotedblleft{}to push\textquotedblright{}, Old High German \emph{stōzan}, Proto-Germanic \emph{stautaną} (EtymDB).
\end{cousinweb}

\begin{grammarlens}[title={What you've built}]
One more word for telling, explaining and describing at B1.
\end{grammarlens}

\subsection*{Your turn}
\begin{itemize}
\raggedright
  \item \emph{Say it:} \emph{der Verstoß}
  \item \emph{Say it:} \emph{der Verstoß}, once more
  \item \emph{Say it:} a sentence of your own with \emph{der Verstoß}
  \item \emph{Recall:} say the German for punishable, then the German for to break, then say \emph{der Verstoß} again
\end{itemize}

\begin{tcolorbox}[breakable,colback=teal!4,colframe=teal!35!black,title={Before you move on}]
What does der Verstoß mean? (\textquotedblleft{}The violation, the breach\textquotedblright{}.) Say it once more.
\end{tcolorbox}
